\documentclass{svproc}
\usepackage{graphicx}%
\usepackage{multirow}%
\usepackage{amsmath,amssymb,amsfonts}%
\usepackage{mathrsfs}%
\usepackage[title]{appendix}%
\usepackage{xcolor}%
\usepackage{textcomp}%
\usepackage{manyfoot}%
\usepackage{booktabs}%
\usepackage{array}%
\usepackage{algorithm}%
\usepackage{algorithmicx}%
\usepackage{algpseudocode}%
\usepackage{listings}%
\usepackage{url}
\usepackage{placeins}
\usepackage{microtype}
\def\UrlFont{\rmfamily}
\newcommand{\Hcal}{\mathcal{H}}
\newcommand{\Ecal}{\mathcal{E}}
\newcommand{\Vcal}{\mathcal{V}}
\newcommand{\R}{\mathbb{R}}
\newcommand{\indicator}{\mathbf{1}}
\newcolumntype{L}[1]{>{\raggedright\arraybackslash}p{#1}}

\begin{document}
\mainmatter

\title{Temporal Directed F-Hypergraphs for Hyperlink Prediction}
\titlerunning{Temporal Directed F-Hypergraphs for Hyperlink Prediction}
\author{Moses Boudourides\inst{1}}
\authorrunning{M. Boudourides}
\institute{School of Professional Studies, Northwestern University, Chicago, IL, USA\\
\email{moses.boudourides@northwestern.edu}}
\maketitle

\begin{abstract}
We formulate temporal directed hyperlink prediction for citation systems represented as singleton-tail directed F-hypergraphs. A citing document is a tail and its contemporaneous bundle of cited documents is a directed hyperedge head. Given a partially observed head at time $t$, the task ranks a hidden head member against historically available candidates. The formulation separates event-level higher-order prediction from ordinary dyadic link prediction and makes temporal admissibility explicit. We define temporal directed F-hypergraph hyperlink prediction, establish the scope of unique-tail reconstruction, and specify a fixed-snapshot chronological evaluation design. The autonomous empirical study treats eight scoring methods on equal grounds: NHP-D-mean, NHP-D-maxmin+, Hyper-SAGNN, CHESHIRE, DHGNN, HPRA, co-citation, and historical in-degree. On identical frozen candidates in Supreme~Court and European Court of Human Rights citation data, no method dominates every corpus and metric. NHP-D-maxmin+ is strongest on the SCOTUS ranking measures, while the smaller ECHR cohort displays metric-dependent results. The paper therefore offers a reproducible comparative benchmark for higher-order prediction on directed citation systems.
\keywords{directed hypergraphs, temporal networks, hyperlink prediction, graph neural networks, citation networks}
\end{abstract}

\section{Introduction}

Many networked systems evolve through events that connect one source to a set of targets: a transaction invokes several instruments, a reaction maps reactants to products, and a document cites a bundle of earlier documents. A simple directed graph represents such an event by independent arcs. This is adequate for many graph-level questions, but it does not make the event head an explicit mathematical object on which a prediction task can be conditioned. Higher-order network models address this issue by representing interactions as hyperedges \cite{berge1973graphs,gallo1993directed}. The directed case is particularly important when source and target roles have distinct semantics.

Citation data provide a tractable real-world instance. Each opinion is a source node and its contemporaneous set of cited precedents is a directed head. The resulting object is a directed F-hypergraph, in the terminology of Gallo et al.~\cite{gallo1993directed}. This paper focuses neither on legal interpretation nor on descriptive jurisprudence. It asks a network-science question: \emph{can a temporally admissible directed-hypergraph model recover a hidden member of a citation-event head?}

The question is a higher-order analogue of link prediction. The prediction task now is conditioned on the remaining members of a set-valued head, candidates must predate the event, and graph-derived inputs cannot use later events. These requirements matter in any temporal directed-hypergraph application, not only in citation networks.

The contribution is fourfold. We (i) define temporal directed F-hypergraph hyperlink prediction and clarify the unique-tail case; (ii) specify a non-leaking chronological candidate-ranking design; (iii) give a common scoring interface and comparable descriptions for six published hypergraph-prediction methods plus two transparent baselines; and (iv) evaluate every method on the same frozen candidates, train--validation--test windows, and calibrated binary decision rule. The paper is deliberately narrow: it offers a mathematical task definition and an autonomous comparative benchmark rather than a theory of judicial institutions.

\section{Temporal Directed F-Hypergraphs}
\label{sec:model}

\begin{definition}[Temporal directed F-hypergraph]
Following \cite{gallo1993directed}, a directed hyperarc is an ordered pair $(T,H)$ of disjoint tail and head subsets of $\Vcal$. An \emph{F-arc} has a singleton tail $T=\{u\}$ and is written $(u,H)$. A temporal directed F-hypergraph is
\begin{equation}
\Hcal=(\Vcal,\Ecal,\tau),\qquad
\Ecal\subseteq\{(u,H):u\in\Vcal,\;\emptyset\neq H\subseteq\Vcal\setminus\{u\}\},
\label{eq:fhypergraph}
\end{equation}
where $\tau:\Vcal\cup\Ecal\to\R$ assigns times to vertices and F-arcs. An F-arc $e=(u,H)$ is \emph{time-respecting} when
\begin{equation}
\tau(v)<\tau(e)=\tau(u),\qquad \forall v\in H.
\label{eq:temporal}
\end{equation}
\end{definition}

For a citation event, $u$ is the citing document and $H$ is its set of cited precedents. The historical state immediately before time $t$ is $\Hcal_{<t}=(\Vcal_{<t},\Ecal_{<t})$, where $\Vcal_{<t}=\{v:\tau(v)<t\}$ and $\Ecal_{<t}=\{e:\tau(e)<t\}$. Thus, condition~\eqref{eq:temporal} rules out a precedent that post-dates its citing opinion.

\subsection{Dyadic projection and the unique-tail case}

The directed dyadic projection is
\begin{equation}
\Pi(\Hcal)=\{(u,v):(u,H)\in\Ecal,\;v\in H\}.
\label{eq:projection}
\end{equation}
The projection is informative, but its information-preservation properties depend on the event structure.

\begin{proposition}[Unique-tail reconstruction]
Let $\mathfrak{F}_1$ be the class of directed F-hypergraphs in which each tail occurs in at most one F-arc. Then $\Pi$ is injective on $\mathfrak{F}_1$.
\end{proposition}
\begin{proof}
For each tail $u$, define $H_u=\{v:(u,v)\in\Pi(\Hcal)\}$. The unique-tail assumption implies that either $H_u=\emptyset$ or the sole F-arc with tail $u$ is $(u,H_u)$. Hence $\Ecal$ is recovered from $\Pi(\Hcal)$. \qed
\end{proof}

The present citation construction deliberately assigns each document one citation-list F-arc, so the premise of Proposition~1 holds. This resolves the apparent tension in the citation setting: the F-hypergraph is not claimed to contain extra raw information relative to an adjacency list with unique tails. Its role is to preserve an \emph{event-level prediction task}: a candidate is scored together with the remaining cited head, not as an isolated dyadic arc. The non-injective case arises in applications in which the same source emits several events, such as recurring transactions, document revisions, or reactions; then the projection loses the partition into event heads.

\section{Time-Respecting Hyperlink Prediction}
\label{sec:task}

For an observed F-arc $e_t=(u_t,H_t)$, choose $v^\star\in H_t$ and set $H_t^-=H_t\setminus\{v^\star\}$. Let $C_t\subseteq\Vcal_{<t}\setminus H_t^-$ be a temporally admissible candidate set that contains $v^\star$. Temporal directed hyperlink prediction ranks
\begin{equation}
s\bigl(u_t,H_t^-\cup\{v\};\Hcal_{<t}\bigr),\qquad v\in C_t.
\label{eq:score}
\end{equation}
The task evaluates whether historical higher-order structure distinguishes the hidden member within a controlled candidate set; it does not ask a model to generate an unrestricted future citation list.

Figure~\ref{fig:formal} depicts the object and chronological evaluation design. A fixed historical snapshot is used: every graph input, random-walk embedding, and candidate popularity statistic is estimated before the train, validation, and test windows. Consequently, later F-arcs cannot alter a ranking at time $t$.

\begin{figure}[t]
\centering
\includegraphics[width=\textwidth]{fig_temporal_farc_model.pdf}
\caption{Mathematical formulation and temporal evaluation design. (a) A directed F-arc has a singleton tail $u_t$ and a set-valued head $H_t$; the prediction task hides $v^\star$. (b) The historical F-hypergraph snapshot precedes the train, validation, and test windows.}
\label{fig:formal}
\end{figure}

\subsection{Candidates, split, and ranking}

Let $B_t=\{v\in\Vcal:\tau(v)<t\}$ be the historically available vertex set. A valid candidate rule obeys
\begin{equation}
C_t\subseteq B_t\setminus H_t^-,\qquad v^\star\in C_t.
\label{eq:admissible}
\end{equation}
This excludes unavailable documents and prevents another positive head member from entering the negative pool. In every experiment, $C_t$ has 20 candidates: the hidden true member, nine non-head nodes with highest historical in-degree, and ten uniformly sampled historical non-head nodes. The mixture creates both popularity-based hard alternatives and ordinary historical alternatives.

After chronological ordering, $q_{.20}$, $q_{.40}$, $q_{.60}$, and $q_{.80}$ denote the 20\%, 40\%, 60\%, and 80\% indices of the strictly time-valid F-arc sequence. The graph snapshot and all candidate statistics use $\Hcal_{<q_{.20}}$ only. Events in $[q_{.20},q_{.40})$ fit model parameters, those in $[q_{.40},q_{.60})$ select the model, and those in $[q_{.60},q_{.80})$ are held out. For each test event,
\begin{equation}
r_t=1+\sum_{v\in C_t\setminus\{v^\star\}}\indicator\!\left\{s(u_t,H_t^-\cup\{v\};\Hcal_{<t})\geq s(u_t,H_t;\Hcal_{<t})\right\}.
\label{eq:rank}
\end{equation}
We report Recall@$k$ (the proportion of hidden members ranked in the top $k$), MRR (the average reciprocal rank), mean rank (the average hidden-member position; lower is better), and candidate-row ROC-AUC. We also report F1, precision, recall, accuracy, log-loss, and Matthews correlation coefficient (MCC) after score calibration and threshold selection on the validation window only. Ties are resolved deterministically by candidate identifier for every method.

\begin{algorithm}[t]
\caption{Fixed-snapshot temporal F-hyperlink ranking}
\label{alg:protocol}
\begin{algorithmic}[1]
\Require Time-respecting F-arcs $\Ecal$; temporal cut points $q_{.20},q_{.40},q_{.60},q_{.80}$
\State Build $\Hcal_{<q_{.20}}$ and all historical candidate statistics
\State Initialize node features by Node2Vec on the snapshot clique expansion
\For{each event $e_t=(u_t,H_t)$ in the later windows}
  \State deterministically hide $v^\star\in H_t$ and form $H_t^-=H_t\setminus\{v^\star\}$
  \State form $C_t=\{v^\star\}\cup C_t^{\mathrm{hard}}\cup C_t^{\mathrm{rand}}$, with $|C_t|=20$
  \State score every candidate F-arc $(u_t,H_t^-\cup\{v\})$, $v\in C_t$
\EndFor
\State fit on $[q_{.20},q_{.40})$, select on $[q_{.40},q_{.60})$, and report $[q_{.60},q_{.80})$
\end{algorithmic}
\end{algorithm}

\section{Models and Comparative Design}
\label{sec:models}

\subsection{Candidate-scoring families}

The study evaluates eight methods through one common prediction task $s_m(e_t,v)$, rather than treating any method as a reference. Two variants of the Neural Hypergraph Predictor for directed hypergraphs (NHP-D) are included: NHP-D-mean and NHP-D-maxmin+ \cite{yadati2020nhp}. For each candidate $v\in C_t$, NHP-D receives $(u_t,H_t^-\cup\{v\})$. Node2Vec \cite{grover2016node2vec} initializes vertices from the fixed historical clique expansion, a graph convolution produces node states, and directed tail--head pooling produces hyperlink and direction logits.

Let $X\in\R^{|\Vcal|\times d}$ be the 64-dimensional Node2Vec matrix and $\widetilde A$ the fixed historical clique-expansion adjacency with self-loops. The NHP-D node states are
\begin{equation}
Z=\sigma\!\left(\widetilde D^{-1/2}\widetilde A\widetilde D^{-1/2}XW\right),
\label{eq:gcn}
\end{equation}
where $\widetilde D$ is the degree matrix \cite{kipf2017semi}. For a candidate F-arc $(u,H)$, with $\bar z_H=|H|^{-1}\sum_{v\in H}z_v$, the singleton-tail local updates are
\begin{equation}
\ell_u=\rho(W_s z_u+W_n\bar z_H),\qquad
\ell_v=\rho\!\left(W_s z_v+W_n\frac{z_u+\sum_{w\in H\setminus\{v\}}z_w}{|H|}\right).
\label{eq:local}
\end{equation}
NHP-D-mean uses $\psi_{\mathrm{mean}}=(\ell_u+|H|^{-1}\sum_{v\in H}\ell_v)/2$, while NHP-D-maxmin+ uses $\psi_{\mathrm{maxmin+}}=\rho(\ell_u-\min_{v\in H}\ell_v)$. Both optimize a hyperlink loss plus a weighted direction loss,
\begin{equation}
\mathcal{L}=\operatorname{BCE}(y,\operatorname{logit}_{\mathrm{hyper}})+0.25\,\operatorname{BCE}(y,\operatorname{logit}_{\mathrm{dir}}).
\label{eq:nhpd-loss}
\end{equation}

Table~\ref{tab:methods} describes each method at the same level: the historical structure it uses and the candidate-head score it supplies. The first six are published prediction families. HPRA~\cite{kumar2020hpra}, co-citation, and historical in-degree are not weaker versions of the neural models; they provide transparent structural comparisons that answer distinct questions about the value added by nonlinear or directed representation learning.

\begin{table}[t]
\caption{Methods evaluated on the common evaluation design. Each receives the same frozen event, observed head, and candidate list.}
\label{tab:methods}
\centering
\footnotesize
\setlength{\tabcolsep}{2.2pt}
\renewcommand{\arraystretch}{.94}
\begin{tabular}{@{}L{.20\textwidth}L{.32\textwidth}L{.42\textwidth}@{}}
\toprule
Method & Historical representation & Candidate-head score\\
\midrule
NHP-D-mean & Clique-expansion GCN with tail--head direction & Mean-pooled local tail/head states \cite{yadati2020nhp}.\\
NHP-D-maxmin+ & Clique-expansion GCN with tail--head direction & Max--min+ local tail/head pooling \cite{yadati2020nhp}.\\
Hyper-SAGNN & Historical node embeddings & Diagonal-masked self-attention over $H_t^-\cup\{v\}$ \cite{zhang2020hypersagnn}.\\
CHESHIRE & Historical head-incidence features & Order-three Chebyshev convolution and max--min plus norm pooling \cite{chen2023cheshire}.\\
DHGNN & Directed historical F-arc propagation & Directed embeddings and candidate--context dot product \cite{ma2024dhgnn}.\\
HPRA & Historical head projection & Resource-allocation attachment to $H_t^-$ \cite{kumar2020hpra}.\\
Co-citation & Historical head projection & Sum of direct head-association weights.\\
Historical in-degree & Historical head frequencies & $\log(1+\text{head frequency}(v))$.\\
\bottomrule
\end{tabular}
\end{table}

\subsection{Common prediction interface and implementation audit}

For every method $m$, the evaluated quantity is the same conditional candidate score,
\begin{equation}
s_m(e_t,v)=s_m\bigl(u_t,H_t^-\cup\{v\};\Hcal_{<q_{.20}}\bigr).
\label{eq:common-score}
\end{equation}
Thus the methods differ in their structural map from the same fixed historical state to a candidate score, not in their temporal information, positives, negatives, or test events. The candidate F-arc is never inserted into the historical encoder. For HPRA, $A$ is the historical degree-preserving head projection and $D_v$ its vertex-degree matrix; its attachment score is
\begin{equation}
s_{\mathrm{HPRA}}(v\mid H_t^-)=\sum_{w\in H_t^-}\bigl(A+AD_v^{-1}A\bigr)_{vw}.
\label{eq:hpra}
\end{equation}

All neural methods are trained on $[q_{.20},q_{.40})$ and selected by validation ranking on $[q_{.40},q_{.60})$. The NHP-D variants are trained anew under this same rule. Published model definitions and publicly available implementations were used to construct conditional adaptations to the common citation-ranking prediction; their differing native data formats are not allowed to change the frozen event or candidate-set specification. Each resulting score file was checked event by event: it contains one positive and 19 negatives, with candidate identifiers identical across methods in every split.

\section{Data and Experimental Design}
\label{sec:data}

The empirical instances are citation F-hypergraphs derived from Supreme Court of the United States (SCOTUS) opinions and European Court of Human Rights (ECHR) judgments. The SCOTUS corpus covers 1791--2024; ECHR is a 10,000-judgment English-language HUDOC sample covering 1959--2013. Citation networks have been used as longitudinal network data \cite{fowler2007network,lupu2012precedent}; here they are used only as directed higher-order event streams.

\begin{table}[t]
\caption{Corpus and strict temporal-evaluation characteristics. The final cohort is smaller than the full graph because every head member and every candidate must be historically available under the fixed-snapshot evaluation design.}
\label{tab:data}
\centering
\begin{tabular}{lrr}
\toprule
& SCOTUS & ECHR\\
\midrule
Citation F-arcs in supplied graph & 24,837 & 18,912\\
Strictly time-respecting F-arcs with $|H|\geq 2$ & 4,793 & 596\\
Historical snapshot F-arcs & 958 & 119\\
Train / validation / test events & 225 / 100 / 61 & 83 / 42 / 42\\
Held-out candidate rows & 1,220 & 840\\
Candidates per event & 20 & 20\\
\bottomrule
\end{tabular}
\end{table}

One head member is hidden deterministically per event. Event-level uncertainty is summarized by 2,000 bootstrap resamples. For the binary measures, each method's raw validation scores are standardized, calibrated by a one-variable logistic map, and assigned the threshold that maximizes validation F1; the fitted map and threshold are then applied once to the held-out scores. This leaves the rank order unchanged and makes F1, precision, recall, accuracy, log-loss, and MCC comparable despite different native score scales. The strict design makes each score a controlled temporal benchmark estimate, not a population estimate of all citations. The ECHR test cohort is particularly small; differences there must therefore be interpreted with its uncertainty intervals rather than as definitive method rankings.

\subsection{Fitting configuration}

The frozen evaluation design applies a single candidate construction and a single held-out cohort to all methods. NHP-D uses 64-dimensional Node2Vec features (20-step walks, ten walks per node, $p=q=1$), AdamW optimization, and validation-window early stopping. Hyper-SAGNN, CHESHIRE, and DHGNN share the same historical vertices and candidate head. Each uses the common validation rule. Architecture-specific settings, score calibration, and binary thresholds are selected without inspecting test ranks or test labels. Changing either candidate construction or test-event eligibility would change the empirical question.

\begin{table}[t]
\caption{Common evaluation contract. The evaluation design is shared by all methods, including the two projection controls.}
\label{tab:configuration}
\centering
\footnotesize
\setlength{\tabcolsep}{3pt}
\begin{tabular}{@{}p{.29\textwidth}p{.61\textwidth}@{}}
\toprule
Component & Specification\\
\midrule
Historical representation & F-arcs strictly before $q_{.20}$; no later F-arc enters a graph-derived feature.\\
Temporal windows & $[q_{.20},q_{.40})$ train; $[q_{.40},q_{.60})$ validation; $[q_{.60},q_{.80})$ test.\\
Candidate set & One hidden member, nine high-in-degree non-head candidates, and ten random historical non-head candidates.\\
Comparison unit & One directed candidate F-arc $(u_t,H_t^-\cup\{v\})$; 20 candidates per event.\\
Model selection & Validation-window ranking; held-out test window untouched until final scoring.\\
Uncertainty & 2,000 event-level bootstrap resamples for Recall@1 and MRR.\\
\bottomrule
\end{tabular}
\end{table}

\subsection{Calibration and metric interpretation}

The native scores of the eight methods have different ranges and meanings: some are probabilities, others are logits, similarities, or resource-allocation values. A common binary threshold on those raw scales would not be meaningful. Therefore, for each method $m$ and court separately, the validation scores are standardized and mapped to a probability by
\begin{equation}
p_m(y=1\mid s)=\left[1+\exp\left\{-a_m\frac{s-\mu_m}{\sigma_m}-b_m\right\}\right]^{-1},
\label{eq:calibration}
\end{equation}
where $(\mu_m,\sigma_m)$ and $(a_m,b_m)$ are fitted from validation rows only. The binary threshold $\theta_m$ is the validation value that maximizes F1. Equation~\eqref{eq:calibration} and $\theta_m$ are then fixed before the test labels are used.

This calibration does not change a method's within-event rank ordering; Recall@$k$, MRR, mean rank, and ROC-AUC therefore retain their direct ranking interpretation. F1, precision, recall, accuracy, log-loss, and MCC instead assess the calibrated retain-or-reject decision. In the one-positive-to-19-negative candidate design, accuracy alone can be misleading. We consequently report it jointly with F1 and MCC. The F1 definition is
\begin{equation}
\mathrm{F1}=\frac{2\mathrm{TP}}{2\mathrm{TP}+\mathrm{FP}+\mathrm{FN}},
\label{eq:binarymetrics}
\end{equation}
and MCC is the correlation of the binary predictions and labels, using all four confusion-matrix cells.

\section{Results}
\label{sec:results}

The autonomous comparison yields no global winner. On SCOTUS, NHP-D-maxmin+ has the strongest Recall@1 (.410), Recall@5 (.770), MRR (.577), and candidate-row ROC-AUC (.824), whereas CHESHIRE has the strongest Recall@10 (.967). On ECHR, co-citation has the strongest Recall@1 (.214) and MRR (.344), Hyper-SAGNN has the strongest Recall@5 (.548), NHP-D-mean has the strongest Recall@10 (.857, tied with NHP-D-maxmin+ and Hyper-SAGNN), and NHP-D-maxmin+ has the highest ROC-AUC (.714). Table~\ref{tab:rankresults} therefore reports an empirical ordering by corpus and metric rather than assuming that the directed models must prevail.

\begin{table}[t]
\caption{Autonomous chronological ranking evaluation on identical 20-candidate test lists. Boldface marks the best result within a court and column.}
\label{tab:rankresults}
\centering
\footnotesize
\setlength{\tabcolsep}{2.6pt}
\renewcommand{\arraystretch}{.94}
\begin{tabular}{llrrrrr}
\toprule
Court & Method & R@1 & R@5 & R@10 & MRR & AUC\\
\midrule
\multirow{8}{*}{SCOTUS}
& Historical in-degree & .000 & .000 & .230 & .078 & .324\\
& Co-citation & .295 & .377 & .656 & .385 & .665\\
& HPRA & .295 & .393 & .639 & .383 & .665\\
& DHGNN & .328 & .689 & .738 & .480 & .767\\
& Hyper-SAGNN & .344 & .738 & .951 & .530 & .798\\
& CHESHIRE & .361 & .754 & \textbf{.967} & .522 & .820\\
& NHP-D-mean & .311 & .754 & .934 & .505 & .785\\
& NHP-D-maxmin+ & \textbf{.410} & \textbf{.770} & .951 & \textbf{.577} & \textbf{.824}\\
\midrule
\multirow{8}{*}{ECHR}
& Historical in-degree & .000 & .190 & .333 & .128 & .378\\
& Co-citation & \textbf{.214} & .476 & .667 & \textbf{.344} & .689\\
& HPRA & .190 & .476 & .667 & .324 & .643\\
& DHGNN & .095 & .476 & .786 & .273 & .691\\
& Hyper-SAGNN & .119 & \textbf{.548} & \textbf{.857} & .308 & .699\\
& CHESHIRE & .095 & .286 & .643 & .230 & .591\\
& NHP-D-mean & .095 & .524 & \textbf{.857} & .302 & .709\\
& NHP-D-maxmin+ & .167 & .476 & \textbf{.857} & .323 & \textbf{.714}\\
\bottomrule
\end{tabular}
\end{table}

Table~\ref{tab:classification} gives the requested binary evaluation. The validation-fitted thresholds test whether a candidate should be retained; this differs from its relative rank within an event. On SCOTUS, HPRA and co-citation have the highest F1 (.409) and MCC (.409), while NHP-D-maxmin+ has the smallest log-loss (.167). On ECHR, NHP-D-mean has the highest F1 (.161), while co-citation has the strongest precision (.308), accuracy (.944), and MCC (.148). These distinctions reinforce the need to report several measures.

\begin{table}[t]
\caption{Validation-calibrated binary test metrics on the same candidate rows. Thresholds maximize validation F1; lower log-loss is better.}
\label{tab:classification}
\centering
\scriptsize
\setlength{\tabcolsep}{1.45pt}
\renewcommand{\arraystretch}{.89}
\begin{tabular}{llrrrrrrr}
\toprule
Court & Method & F1 & Prec. & Rec. & Acc. & AUC & Loss & MCC\\
\midrule
\multirow{8}{*}{SCOTUS}
& Historical in-degree & .166 & .091 & \textbf{1.000} & .499 & .676 & .186 & .207\\
& Co-citation & \textbf{.409} & .594 & .311 & \textbf{.955} & .665 & .177 & \textbf{.409}\\
& HPRA & \textbf{.409} & .594 & .311 & \textbf{.955} & .665 & .177 & \textbf{.409}\\
& DHGNN & .263 & \textbf{.667} & .164 & .954 & .767 & .177 & .316\\
& Hyper-SAGNN & .330 & .405 & .279 & .943 & .798 & .172 & .307\\
& CHESHIRE & .299 & .500 & .213 & .950 & .820 & .174 & .305\\
& NHP-D-mean & .291 & .327 & .262 & .936 & .785 & .175 & .260\\
& NHP-D-maxmin+ & .265 & .500 & .180 & .950 & \textbf{.824} & \textbf{.167} & .280\\
\midrule
\multirow{8}{*}{ECHR}
& Historical in-degree & .127 & .070 & \textbf{.762} & .479 & .622 & .194 & .099\\
& Co-citation & .145 & \textbf{.308} & .095 & \textbf{.944} & .689 & .192 & \textbf{.148}\\
& HPRA & .143 & .286 & .095 & .943 & .643 & .193 & .141\\
& DHGNN & .151 & .086 & .619 & .651 & .691 & .188 & .123\\
& Hyper-SAGNN & .132 & .088 & .262 & .827 & .699 & .190 & .073\\
& CHESHIRE & .062 & .055 & .071 & .892 & .591 & .196 & .006\\
& NHP-D-mean & \textbf{.161} & .094 & .571 & .702 & .709 & .185 & .133\\
& NHP-D-maxmin+ & .158 & .099 & .381 & .796 & \textbf{.714} & \textbf{.184} & .110\\
\bottomrule
\end{tabular}
\end{table}

Figure~\ref{fig:comparators} displays MRR with 95\% event-bootstrap intervals. It visualizes the principal ranking result without privileging a method a priori: SCOTUS favors NHP-D-maxmin+ on MRR, whereas the ECHR intervals are broad and overlap substantially. In conjunction with Tables~\ref{tab:rankresults}--\ref{tab:classification}, this is evidence for context-dependent comparative performance, not universal superiority.

\begin{figure}[t]
\centering
\includegraphics[width=.82\textwidth]{fig_comparator_mrr.pdf}
\caption{Held-out MRR with 95\% event-bootstrap intervals for all methods. Every score ranks the same hidden member against the same 19 negative candidates per event.}
\label{fig:comparators}
\end{figure}

\FloatBarrier

\section{Validity Conditions and Limitations}
\label{sec:validity}

The benchmark measures conditional ranking, not unrestricted generation of a future head. The partial head $H_t^-$ is observed by construction and the model selects the hidden member from a finite historical candidate set. The design is valid only when three conditions hold: temporal validity requires~\eqref{eq:temporal}; representation validity requires that graph-derived inputs use only the fixed historical snapshot; and event validity requires that candidate $v$ be evaluated in the displayed F-arc $(u_t,H_t^-\cup\{v\})$ rather than as a projected pair in isolation.

The comparison has corresponding limits. Strict temporal eligibility leaves 61 SCOTUS and 42 ECHR test events, and the 20-candidate pool is sampled rather than exhaustive. Models with a head-only candidate architecture are evaluated on the displayed head; NHP-D and DHGNN additionally use directed historical roles. The implementations are documented conditional adaptations of published architectures to one citation-ranking prediction. A rolling snapshot could instead update representations after each event, but would introduce online retraining and candidate-refresh choices outside the present controlled comparison.

\section{Discussion and Conclusion}

Temporal directed F-hypergraph hyperlink prediction is well-defined only after the event boundary, historical candidate universe, hiding rule, and time-respecting feature construction are specified. The unique-tail proposition clarifies that a citation-list F-hypergraph need not add raw information beyond a complete adjacency list; its contribution is to preserve the set-valued event on which the prediction question is conditioned.

The comparative results establish a method-dependent empirical ordering rather than a preassigned winner. On SCOTUS, NHP-D-maxmin+ leads the principal ranking measures, while CHESHIRE leads Recall@10 and the projection methods lead some calibrated binary measures. On ECHR, co-citation leads Recall@1 and MRR, Hyper-SAGNN leads Recall@5, NHP-D-mean shares the best Recall@10 and has the best F1, and NHP-D-maxmin+ has the best ROC-AUC and log-loss. Thus, the contribution is an autonomous, reproducible temporal task and a comparative result showing how directed, set-aware, spectral, and projected representations behave on the same event-level question.

\bibliographystyle{spmpsci}
\bibliography{references}
\end{document}
